% A.63: full independent root and DE proof/source/certificate audits PASS.
\subsection{Every isogeny enlargement of an ordinary Heisenberg cover}

\begin{statement}{A.63}{Proved}{All ordinary Heisenberg enlargements have cusp fixed groups and normalizers}
Let \(p\equiv1\pmod4\) be prime. Let \(C_p\) be the smooth
proper normalization of \(z^p=\psi_p(x)\) over
\(E:y^2=x^3-x\), with its actual degree-\(p\) map \(q\).
For every integer \(m\ge1\), put
\[
 S_{p,m}=q^{-1}(E[pm]),\qquad
 Y_{p,m}=C_p\setminus S_{p,m}
       \xrightarrow{[pm]q} E\setminus\{O\}.
                                                        \tag{A.63.1}
\]
This is a finite etale cover of degree \(p^3m^2\), with
\[
 g(C_p)=1+\frac{p^2(p-1)}2,\qquad
 |S_{p,m}|=p^3m^2-p^3+p^2.                              \tag{A.63.2}
\]
Both assertions of Task A.5 hold for every
coefficient prime, every selected cusp packet, and every
permitted open derived subgroup for these actual covers.

More precisely, at a place above \(p\), every cusp has a
distinct cusp with the same specialization on the unmarked
proper stable model. In the stable marked model their joining
path contains only rational components. No assertion that
odd-order torsion avoids the positive-genus tails is used.
There is no assertion that this family is cofinal among all
finite covers.
\end{statement}

\begin{proof}
We work after a finite number-field extension making the
actual deck group and the finite marking set rational, and
after a further extension at a place above \(p\) admitting
the proper stable model. All valuations below are normalized
by \(v(p)=1\). Such finite extensions are permitted in the
open-derived comparison of A.41--A.42.

\paragraph{The actual cover and its finite deck action.}
By A.62, the degree-\(p^3\) map \([p]q\) has the actual
exponent-\(p\) Heisenberg group \(G\). Its centre \(Z\)
has order \(p\), is generated by
\(\sigma:z\mapsto\zeta_pz\), and fits in
\[
 1\longrightarrow Z\longrightarrow G
       \longrightarrow E[p]\longrightarrow1.
                                                        \tag{A.63.3}
\]
The last map records the actual translation on the target of
\(q\). The only nontrivial stabilizers of geometric points
of \(C_p\) under \(G\) are exactly \(Z\), at the original
cusps \(q^{-1}(E[p])\). Indeed a fixed point under a deck
element must map under \(q\) to a point fixed by its target
translation; a nonzero translation has no fixed point.
The residual possibility is the central cyclic cover \(q\),
whose branch locus is exactly \(E[p]\), with index \(p\)
at every point. This follows also from
\(\operatorname{div}(\psi_p)=D_{E[p]}-p^2(O)\).

The marking set \(S_{p,m}\) is \(G\)-stable, since
translations by \(E[p]\) preserve \(E[pm]\).
The composite in (A.63.1) is etale away from this set:
\([pm]\) is etale in characteristic zero, and the entire
branch set of \(q\) has been removed. Its degree is
\(p(pm)^2\). There are \(p^2\) branch fibres with one
point each, and \(p^2m^2-p^2\) other fibres with \(p\)
points each. This proves (A.63.2), including the exact cusp
count and the Euler identity \(2g-2+|S_{p,m}|=p^3m^2\).

\paragraph{The small connected translations.}
Let \(E[p]^0\subset E[p](\overline K_v)\) be the subgroup
specializing to the origin of the ordinary good reduction.
It has order \(p\). For a formal parameter \(t\) at \(O\),
every nonzero \(T\in E[p]^0\) satisfies
\[
                         v(t(T))=\frac1{p-1}.           \tag{A.63.4}
\]
Here is the exact valuation argument. The formal group is
defined over \(\Z_p\), and ordinarity makes its reduction
have height one. Thus
\[
 [p](t)=pt+\sum_{2\le j<p}p b_jt^j+u_pt^p
                       +\sum_{j>p}u_jt^j,
 \quad b_j,u_j\in\Z_p,\quad u_p\in\Z_p^\times.
                                                        \tag{A.63.5}
\]
The divisibility of the coefficients with \(2\le j<p\)
is \cite[Corollary IV.4.4, p.126]{fghSilvermanAEC};
the height and inseparable-degree comparison is
\cite[Theorem IV.7.4, p.134]{fghSilvermanAEC}.
For a nonzero root with \(v(t)>0\), the terms with
\(2\le j<p\) have valuation strictly larger than
\(1+v(t)\), and those with \(j>p\) have valuation
strictly larger than \(p v(t)\). The two remaining terms
must have equal valuation. This proves (A.63.4).

Define the actual order-\(p^2\) subgroup
\[
             G_c=\operatorname{pr}^{-1}(E[p]^0)\subset G.
                                                        \tag{A.63.6}
\]
It contains \(Z\) and is isomorphic to
\((\Z/p)^2\): every group of order \(p^2\) is abelian,
and the ambient group has exponent \(p\).

\paragraph{Its action on every positive-genus component.}
Use the actual proper stable fibre constructed in A.62.
It has one elliptic main component, mapping radicially of
degree \(p\) to the good elliptic special fibre, and the
following other positive-genus components:
\[
 V^p-V=U^k,\qquad k=2\text{ or }4.                       \tag{A.63.7}
\]
There is one such tail above each critical centre specified
in A.62. The genus exhaustion there proves that these are
all positive-genus components and that the dual graph is a
tree. Any remaining proper-stable component is rational.
The group action extends to this stable model by
\cite[Proposition 10.1.16 and Theorem 10.3.34(b),
pp.460--461 and 518--519]{fghLiu}.

Every element of \(G_c\) reduces to the identity translation
on the base elliptic curve. Hence it fixes the main component
and acts identically on its function field: the inclusion
from the base is purely inseparable, which admits no
nontrivial automorphism over the base.

We check the tail action in its actual scaled chart. Let
\(s\) be the good integral parameter at its unramified
centre used in A.62: it is \(x-a\) away from two-torsion,
and \(y\) at the relevant two-torsion centre. The tail
valuation is represented by
\[
 s=\rho U,\qquad
 v(\rho)=\frac1{k(p-1)},\qquad
 z=g z_a(1+\eta V),\quad v(\eta)=\frac1{p-1}.
                                                        \tag{A.63.8}
\]
In the two-torsion chart \(x-a\) has order \(s^2\), as
in A.62; the parameter in (A.63.8) is still \(s=y\).
Any constant rescaling making (A.63.7) monic is a unit and
does not change these valuations.

The smooth integral group law gives, on the completed
neighbourhood of this centre,
\[
                 s(P+T)-s(P)\in t(T)R[[s,t(T)]].        \tag{A.63.9}
\]
To justify this inclusion, write translation by a formal
point with parameter \(t\) as an integral power series
in \(s,t\). At \(t=0\) it is the identity, so the
difference is divisible by \(t\). All coefficients are
integral because the centre is a smooth point of the good
elliptic group scheme. Evaluating the tail Gauss valuation
and using (A.63.4) gives
\[
 v\bigl(s(P+T)-s(P)\bigr)\ge\frac1{p-1}
           >\frac1{k(p-1)}=v(\rho).                    \tag{A.63.10}
\]
Thus translation by \(T\in E[p]^0\) preserves the actual
tail disc and induces the identity on its reduced variable
\(U\). There is exactly one positive-genus component
above this centre, so every lift in \(G_c\) preserves the
tail. Its action on (A.63.7) fixes \(U\); consequently it
is \(V\mapsto V+c\) for some \(c\in\F_p\), because
these are exactly the automorphisms of the degree-\(p\)
Artin--Schreier extension over \(\overline k(U)\).

The central generator acts nontrivially on this tail. In
the actual chart (A.63.8) it sends
\[
       V\longmapsto\zeta_pV+\frac{\zeta_p-1}{\eta}.
                                                        \tag{A.63.11}
\]
The second summand is a unit, since
\(v(\zeta_p-1)=v(\eta)=1/(p-1)\); its reduction is a
nonzero element of \(\F_p\), as also follows by
preserving (A.63.7). The tail-action homomorphism is
therefore a surjection
\[
 G_c\longrightarrow\F_p,\qquad
 K_B:=\ker(G_c\longrightarrow\operatorname{Aut}(B)),
 \quad |K_B|=p,\quad K_B\cap Z=1.                       \tag{A.63.12}
\]
Here \(B\) denotes that individual tail. In particular
\(K_B\) is a noncentral subgroup acting identically on
the whole smooth proper tail, including its point at
infinity. This is an actual group kernel, not a kernel of
an approximate arithmetic representation.

\paragraph{Rational components and the exact marked specialization.}
Every positive-genus vertex of the proper stable tree is
fixed by \(G_c\), by the preceding analysis. Every leaf
of that tree is positive-genus: an unmarked rational leaf
would be unstable. An automorphism of a finite tree fixing
all its leaves fixes the whole tree. For example, every
non-leaf vertex is determined by the paths between the
leaves in its incident branches; vertices of valence two
also occur on such fixed paths. Consequently \(G_c\)
fixes all vertices and edges. Each rational proper-stable
component has at least three distinct nodes, and these
nodes are fixed individually. An automorphism of
\(\mathbf P^1\) fixing three points is the identity.
Thus \(G_c\) acts pointwise on every rational component
and fixes every node of the proper stable fibre.

Let \(P\in S_{p,m}\), and let \(x\) be its specialization
on this unmarked stable model. If \(x\) is a smooth point
of an Artin--Schreier tail \(B\), choose a nonidentity
element of \(K_B\). If \(x\) is on the main component,
on a rational component, or is a node, choose any element
of \(G_c\setminus Z\). In either case the element fixes
\(x\) and does not fix \(P\), by the exact generic
stabilizer statement following (A.63.3). We obtain
\[
                    Q\ne P,\qquad Q\in S_{p,m},
                    \qquad \operatorname{sp}(Q)=x.
                                                        \tag{A.63.13}
\]
In fact the chosen cyclic subgroup gives \(p\) distinct
marked points with that specialization. The choice may
depend on the component; no single noncentral subgroup is
asserted to act identically on every tail at once.

We spell out why (A.63.13) is the pointed hypothesis
needed in A.41--A.42. Forget the marks in a semistable
marked model and stabilize. This contracts only rational
components. Above a smooth point of the unmarked model
the exceptional fibre is a rational tree attached at that
point, and above a node it is a rational chain with
possibly attached rational trees. The fibres are
connected, and the ambient dual graph is a tree. Since
the marked sections are disjoint and meet only smooth
points, two distinct sections specializing to the same
unmarked point are separated on this rational exceptional
fibre. Neither can stay at an attaching node of a
positive-genus component. Their joining path, including
both endpoints, therefore contains only rational
components. The same conclusion is immediate if their
unmarked component is already rational. This describes
the contraction fibre, not just the image point on the
elliptic base.

\paragraph{The actual fixed group and normalizer conclusions.}
The rational companion just constructed works for every
selected cusp. A.41--A.42 apply directly to \(Y_{p,m}\),
at residue characteristic \(p\), and give
\[
 \bigl(\pi_1(Y_{p,m,\overline L})^{(\ell)}\bigr)^{s(\Delta_0)}
                        =I_{P,\ell}                    \tag{A.63.14}
\]
for every coefficient prime \(\ell\), every permitted
packet, and every permitted open derived subgroup.
As there, choose a finite field of definition whose
closed derived subgroup lies in the given \(\Delta_0\);
the latter already fixes the selected inertia, so the
conclusion passes to it. This is a direct geometric
application, not a fixed-group lift through a quotient.

For normalizers use the actual open immersion
\(Y_{p,m}\hookrightarrow Y_{p,1}\). It induces the
identity on the proper Jacobian. If a noncuspidal closed
procyclic subgroup of the geometric pro-\(\ell\) group
had open arithmetic normalizer, A.11 would give it
nonzero image in proper rational homology. Its image in
the target pro-\(\ell\) group is then infinite
procyclic, still has open arithmetic normalizer, and has
the same nonzero proper-homology image. A.62 puts this
image in target cusp inertia, which has zero proper
homology, a contradiction. No assertion about the
nilpotent Galois closure of the enlarged composite is
needed.

The exact certificate \path{code/a_heisenberg_ordinary_enlargements.py}
checks the counting identities, both strict radius
inequalities, all tail-kernel characters in several
actual finite Heisenberg models, and an ordinary
critical point of exact odd order. The uniform argument
above proves the statement for all \(p,m\).
\end{proof}
